% Français 9115, batch 36 — Gallica views 701–720.
% Pass: Opus 5.5 (claude-opus-5-5), 2026-09-22 — first pass, unchecked against the views by a human.
%
% All twenty views are on disk in the mirror; none is missing. Batches 31–35
% were being transcribed in parallel and were not read; view 700 and view
% 721 were looked at as thumbnails only, to see what crosses the edges.
%
% What the twenty views hold. Loose leaves and slips of her own, each
% mounted low on a grey mounting sheet, recto and verso photographed; most
% are number theory, and among them are the three pages that Laubenbacher
% and Pengelley cite as « Manuscript C », ff. 348r–349r. Twelve views carry
% text; eight are blank backs (v. 702, 703, 705, 707, 711, 714, 715, 717),
% whose reversed writing is the other side read through the paper, checked
% by mirroring.
%
%   v. 701      Written side of a small leaf whose other side (v. 700, not
%               in this batch) carries the ink number 344 and is blank.
%               Heat on a hemisphere: the surface divided by meridians and
%               parallels into infinitely small quadrilaterals, a heat
%               source on a meridian side of a parallelogram resting on the
%               equator, the temperatures f f′ f″ … and a a′ a″ … of the
%               first two rows. A draft, stopping mid-sentence a third of
%               the way down (« … sur le premier parallèle »).
%   v. 704      Leaf 346. « Voici encore une autre demonstration dela
%               propriété du nombre 2 d'être résidu ou non résidu des
%               nombres premiers suivant leurs formes par rapport a 8 »:
%               the cases 4K+3 (K even / odd) and 4n+1 with n divisible by
%               an odd power of 2, by −1 ≡ 2(4K′+1) etc.; ends on a dash,
%               « cette methode n'apprend rien pour le cas p = 2^{2m}n+1 ».
%               Complete on its leaf.
%   v. 706      Leaf 347, headed « Observations sur les nombres premiers
%               de la forme 4n+1 considérés comme résidus d'autres nombres
%               premiers demême forme »: X² + P = Mp written with P, p sums
%               of two squares and M of four, and the ways of equating the
%               terms; a « N.ota » and a margin note on products of numbers
%               decomposable into four squares in only one way. A « 1/ » at
%               the top left. Breaks off at the foot of the leaf on an
%               underlined « 2 »; the back (v. 707) is blank and the next
%               written view is another text.
%   v. 708–710  MANUSCRIPT C (ink 348 on v. 708, its back v. 709, ink 349
%               on v. 710). A fair copy in a regular, posed hand, almost
%               without correction, headed « Théorème » / « Démonstration »:
%               for n odd > 1, 2z^{2n} = y^{2n} + x^{2n} is impossible in
%               integers (the leaf writes z where x is meant in the
%               statement, x four lines later), proved by writing z^n and
%               the other terms as sums and differences of squares
%               (systems (1) and (2)), then the divisibility of
%               2(v²−u²)vu by pq(p²−q²) and the two alternatives
%               2(v²−u²)vu > pq(p²−q²) and = 0, down to u = 1, v = 0 and
%               n = 1; ends at v. 709 on « Ce cas est donc le seul … ».
%               On v. 710: the same proof extended to 2z^{2n} =
%               y^{2(n−2d)} + x^{2(n−2d)}, hence a theorem for 2z^{2(n+2d)}
%               = x^{2n} + y^{2n}, n odd; then a second « Théorème »: « Pour
%               toute valeur de n plus grande que l'unité il est impossible
%               de satisfaire en nombres entiers a l'équation Z^{2n} =
%               X^{2n} + Y^{2n} », with its « Démonstration » reducing it,
%               by X = 2xx′ and Z^n ± Y^n, to 2(x)^{2n} = (x)′^{2n} +
%               (x)″^{2n}, « équation dont l'impossibilité vient d'être
%               demontrée pour toute valeur de n plus grande que l'unité ».
%               Three written pages, 348 recto, 348 verso, 349 recto, one
%               text from the first word to the last; the back of 349
%               (v. 711) is blank. These are the views and numbers of
%               ff. 348r–349r.
%   v. 712–713  Leaf 350, both sides, a fast hand: heat again. The
%               Laplacian of v in the coordinates r, θ, φ, term by term
%               (d²v/dy², d²v/dx² + d²v/dy², d²v/dz²), a « N.ota » on the
%               differential of cos θ, the sum; r infinite gives the
%               cylinder, « celle que M.r Poisson a donné en considérant le
%               cas général », and with no flow along the axis « celle que
%               j'avois trouvé »: the circular plate, « abstraction faite de
%               la déperdition de chaleur due dans la plaque isolée au
%               contact de l'air »; θ = 90; then on the back how the
%               equation becomes that of the plate, and « On peut arriver a
%               l'équation du cône » by the heat entering the three faces of
%               the element, K r² sin θ cos θ dθ dφ · dv/dr … Stops three
%               quarters down v. 713.
%   v. 716      Leaf 352, headed « Cone »: the full working of the same
%               change of variables (dr/dx, d.sin θ/dx, dφ/dx …, d²v/dx²
%               developed, d²v/dy², d²v/dz², their sum), to the cylinder
%               again at the foot. Nothing orders it against leaf 350,
%               which gives the same results more cleanly.
%   v. 718–719  Leaf 353 (top) and a small leaf 354 (below) on one
%               mounting sheet. 353: two lines only on the numbered side
%               (« est non résidu d'un nombre premier de la même forme
%               4k+3 »); its back (v. 719) is written upside down relative
%               to the view and crossed out with a large cross: 120k+1 and
%               40k+11, 2 and 5 as residues, and a table of products ≡ 11
%               (mod 40) with the classes q = 40k+29, 19, 21, 27, 13, 37.
%               354, both sides: « Ici, l'auteur ne dit rien des nombres de
%               la forme 4k+3. J'observe que l'on peut démontrer que les
%               nombres 12n+7 … sont non résidus pour quelques nombres
%               premiers moindres », by 3aa − p and the factors of 12K+5;
%               the back does the same for 12n+5.
%   v. 720      Leaves 355 and 356 on one mounting sheet, the library's
%               round stamp on 355. 355: « En examinant la manière dont
%               Gauss démontre (N.o 125) que tout nombre premier p = 4n+1
%               pris negativement est non résidu d'un nombre premier moindre
%               que p … », with a reference to « la feuille ci-jointe ».
%               356: « Je me propose de démontrer qu'un nombre quelconque
%               de la forme 4K+3, pris, tant avec le signe +, qu'avec le
%               signe −, est non résidu de quelque nombre premier moindre
%               que lui », −p and then p in the cases 8n+3 and 8K+7; a
%               remark at the foot, « Le but de l'auteur étant de passer au
%               théorème fondamental … », breaking off at the leaf's edge
%               on « pour en ». Its back (v. 721, thumbnail only) is blank.
%               The « auteur » of 354 and 356 is not named on those leaves;
%               that it is Gauss, the Disquisitiones, is suggested by 355
%               and is an inference.
%
% The numbers on the leaves, and why no \folio{} is written. (1) The
% top-right ink series, one number per leaf, on the side that carries it:
% 344 (v. 700, out of batch), 345 (v. 702, a blank recto), 346 (v. 704),
% 347 (v. 706), 348 (v. 708), 349 (v. 710), 350 (v. 712), 351 (v. 714,
% blank), 352 (v. 716), 353 and 354 (both on v. 718), 355 and 356 (both on
% v. 720). No gap. The 5 of this series has a long descending tail and
% reads like ζ; the 4 of 354 is closed. On v. 714 a number on the
% projecting top edge of a later leaf reads 3·3, its middle figure hidden:
% the edge is that of leaf 353. The written backs carry no number. The
% series is penned, as everywhere in this volume, so no \folio{} is
% written (hand.md, « The open question of this volume »); but it
% coincides exactly with the folios Laubenbacher and Pengelley cite for
% Manuscript C (348r–349r), which is one more piece of evidence that this
% ink series is the foliation the literature uses. (2) Her own marks: a
% « 1 » followed by an oblique stroke at the top left of v. 706; system
% labels (1), (2) at v. 708–710. (3) Numbers cited: Gauss « N.o 125 »
% (v. 720). Every number was read on its view; none is inferred.
%
% Notation and conventions. Her spelling is kept and not flagged:
% « dequadrilatères », « ala », « aumoins », « dela », « laforme »,
% « parconséquent », « c'est adire », « fesoit », « fesant », « cequi »,
% « cequ'il », « cidessus », « demême », « quelquesnombres », « nedit »,
% « tems », « appuyent », « concluerons », « des ces remarque ». Her
% congruence sign is three bars, set \equiv; her « plus petit que » is a
% small Z on its side, set $<$. sin, cos are written with a long s and set
% \sin, \cos, with her stops kept where they are clear (« sin. θ »,
% « cos. φ »). Her « 4 » is shaped like an L and her « 3 » like a 9 with a
% long tail (hand.md, v. 301, 377); both were checked against the
% arithmetic. Underlined letters in formulas stay \underline; underlined
% words in running text are given as \emph{}. Capitals and lower-case
% letters that the leaf keeps apart (U, V, u, v at v. 708–709; Z, X, Y and
% x, x′ at v. 710; P, X and p, x at v. 706) are kept apart. Where the
% leaf and the calculation disagree the leaf stands, with a note: v. 708
% (z for x in the statement), v. 709 ((v²+u²) for (v²−u²) in two
% inequalities), v. 710 (an exponent read 8n where 2n is expected; 2^{n−1}
% beside 2^{2n−1}), v. 712 (d²u for d²v), v. 713 (d²v/dx), v. 716
% (d²v/dy, a denominator missing).
%
% What could not be read. \ill{} stands in five places: a word after
% « négligeable ala » (v. 701); in the struck first numerator of a brace
% (v. 716); a struck sign or letter after a parenthesis (v. 716); a blotted
% sign after « immédiatement » and a heavily struck word before
% √(p/2) (v. 720). Readings offered with doubt (\uncertain{}): « Si »,
% « seront », « fait » (v. 701); « 8n » and the « 2 » of 2(n+2d) (v. 710);
% φ at the edge and two signs under blots (v. 712); three signs under
% blots (v. 716); « ou 11 », « ou 19 » (v. 719); the second « 4n+1 »
% (v. 720); « aisée », « remplir », « pro. » (v. 720), and « produisent »
% in the margin note of v. 706. Nothing was guessed to fill a gap.
%
% Nothing here was checked against any published edition or study.
\documentclass[11pt,a4paper]{article}
% The preamble shared by every transcription of the mathematicians' manuscripts in Gallica.
%
% It rests on one idea: **uncertainty compounds, it does not resolve.** A
% transcription that smooths over a doubtful word has destroyed the only thing
% it was made for — a reader must be able to tell, at every point, what was on
% the page and what was supplied. The apparatus macros below are therefore the
% critical apparatus, not decoration.
%
% The same commands are recognised by `scripts/render.mjs`, which produces the
% site's reading view, and by `scripts/tei.mjs`. A macro added here must be
% added there too, or rendering fails loudly — which is the wanted behaviour.
%
% Comments here are English, like the rest of the repository. The documents
% this preamble sets are French, and that is deliberate: see the skills.

\ProvidesPackage{germain}[2026/09/15 Transcriptions of mathematicians' manuscripts in Gallica]

\usepackage{amsmath,amssymb,amsthm}
\usepackage{mathrsfs}
% Kept from the parent preamble so the renderer's diagram support stays
% available; these manuscripts draw no commutative diagrams, and nothing here uses it.
\usepackage{tikz-cd}
\usepackage[english,french]{babel}
\usepackage{fontspec}

% --- Greek ------------------------------------------------------------------
% The text face stays fontspec's default, Latin Modern, which has no Greek:
% XeTeX drops every glyph it lacks with a line in the log and nothing on the
% page, and the Greek words the manuscripts quote vanished from the PDFs. So
% the Greek and Greek Extended blocks get a XeTeX character class of their own,
% and entering or leaving a run of it switches the family to CMU Serif — the
% same Computer Modern design, with polytonic Greek — and back. Only the
% family changes: a Greek word in \textit or \textbf keeps its shape and series.
% The transitions to and from every other class carry the tokens that class
% has towards an ordinary letter, so French spacing before « ; » or inside
% « » still holds after a Greek word. No grouping, so no brace or \endgroup
% can come out unbalanced; maths is left alone, since XeTeX inserts nothing
% there. Kept identical in `exercices.sty`.
\makeatletter
\newfontfamily\germainGreekfont{cmun}[
  Extension=.otf, NFSSFamily=germaingreek,
  UprightFont=*rm, ItalicFont=*ti, SlantedFont=*sl,
  BoldFont=*bx, BoldItalicFont=*bi, BoldSlantedFont=*bl]
\newXeTeXintercharclass\germain@greek
\def\germain@greekrange#1#2{%
  \@tempcnta=#1\relax
  \loop
    \XeTeXcharclass\@tempcnta=\germain@greek
    \ifnum\@tempcnta<#2\advance\@tempcnta\@ne\repeat}
\germain@greekrange{"0370}{"03FF}
\germain@greekrange{"1F00}{"1FFF}
\def\germain@greekprev{\rmdefault}
\def\germain@greekin{%
  \edef\germain@greekprev{\f@family}\fontfamily{germaingreek}\selectfont}
\def\germain@greekout{\fontfamily{\germain@greekprev}\selectfont}
\def\germain@greekjoin{%
  \XeTeXinterchartokenstate=\@ne
  \@tempcnta=\z@
  \loop
    \ifnum\@tempcnta=\germain@greek\else
      \XeTeXinterchartoks\germain@greek\@tempcnta=\expandafter{%
        \expandafter\germain@greekout\the\XeTeXinterchartoks\z@\@tempcnta}%
      \XeTeXinterchartoks\@tempcnta\germain@greek=\expandafter{%
        \the\XeTeXinterchartoks\@tempcnta\z@\germain@greekin}%
    \fi
    \ifnum\@tempcnta<4095\advance\@tempcnta\@ne\repeat}
% babel-french sets its own classes, and saves and restores the interchar
% state, each time a language is selected: join again after it does.
\addto\extrasfrench{\germain@greekjoin}
\addto\extrasenglish{\germain@greekjoin}
\AtBeginDocument{\germain@greekjoin}
\makeatother

\usepackage{geometry}
\geometry{a4paper,margin=25mm,marginparwidth=28mm}
\usepackage{marginnote}
\usepackage{xcolor}
\usepackage[normalem]{ulem}
\setlength{\emergencystretch}{3em}
\usepackage{hyperref}
\hypersetup{colorlinks=true,linkcolor=black,urlcolor=black,pdfborder={0 0 0}}

% --- Batch metadata ---------------------------------------------------------
% Taken as they stand by the rendering script, which shows them at the head of
% the reading view. They fix what the file claims to cover. The macro names are
% the parent project's, so that the scripts stay shared; \folder here means the
% volume's slug — `fr-9115` — and \pages counts Gallica views.

\newcommand{\folder}[1]{\gdef\grothFolder{#1}}
\newcommand{\batch}[1]{\gdef\grothBatch{#1}}
\newcommand{\pages}[2]{\gdef\grothFirst{#1}\gdef\grothLast{#2}}
\newcommand{\foldertitle}[1]{\gdef\grothFoldertitle{#1}}

% The holder's own shelfmark, printed as they write it: « Français 9115 ».
\newcommand{\shelfmark}[1]{\gdef\grothShelfmark{#1}}

% The dating, copied verbatim from the holder's catalogue — never inferred
% here. « XIXe siècle » is the BnF's; a pass that dates a leaf from its content
% says so in a \note{}, as its own inference.
\newcommand{\dating}[1]{\gdef\grothDating{#1}}

% The status notice, printed in the title block and repeated as a diagonal
% watermark on every page of the PDF. The manuscripts are in the public
% domain; what this line declares is not a right but a quality — an
% unverified first pass by a machine. The skills write the line; a file
% without it gets no watermark, so leaving it out is visible in the output.
\newcommand{\watermark}[1]{\gdef\grothWatermark{#1}}

\gdef\grothFolder{?}\gdef\grothBatch{}\gdef\grothFirst{?}\gdef\grothLast{?}%
\gdef\grothFoldertitle{}\gdef\grothShelfmark{}\gdef\grothDating{}\gdef\grothWatermark{}

% --- The résumé -------------------------------------------------------------
\newenvironment{resume}
  {\small\setlength{\parskip}{0.45em}}
  {\par\medskip\hrule\medskip}

% --- Keywords ---------------------------------------------------------------
% In English, closing the modernised reading's résumé: the modern vocabulary
% under which to search for what the leaves construct. The manifest extracts
% them to tag the volume — this line is the single source of the tags.
\newcommand{\keywords}[1]{\par\smallskip\noindent
  {\footnotesize\textsc{Keywords} --- \textit{#1}}\par}

% --- The critical apparatus -------------------------------------------------

% The Gallica view number, in the margin. This is the anchor that lets a
% disagreement over a formula be settled by looking at *one* image rather than
% a volume of 750. The site uses it to turn the facsimile as the transcription
% is scrolled. A view is one scanned image: usually one side of a leaf.
\newcommand{\page}[1]{%
  \marginnote{\footnotesize\color{gray!70}\textsf{v.\,#1}}%
  \label{page:#1}\ignorespaces}

% The BnF's pencilled foliation, where the leaf carries it and a pass can read
% it: « 198r ». This is what the literature cites — « ff. 198r–208v » — and
% the one line that turns a citation into a view. Recorded, never inferred:
% a folio a pass cannot read is not written.
\newcommand{\folio}[1]{%
  \marginnote{\footnotesize\color{gray!50}\textsf{f.\,#1}}\ignorespaces}

% The modernised reading's coarser counterpart to \page{}: which views a
% section is drawn from.
\newcommand{\pagerange}[2]{%
  \marginnote{\footnotesize\color{gray!70}\textsf{v.\,#1--#2}}%
  \ignorespaces}

% Illegible: never guessed.
\newcommand{\ill}{\textcolor{red!60!black}{[\,\ldots\,]}}

% A doubtful reading: the word is given, and flagged as uncertain.
\newcommand{\uncertain}[1]{\uline{#1}}

% An editorial addition — a missing letter, an implied word.
\newcommand{\add}[1]{\textcolor{blue!50!black}{[#1]}}

% Struck out by the writer. What was crossed out is part of the document.
\newcommand{\struck}[1]{\sout{#1}}

% The transcriber's note, on the state of the page or the mathematics.
\newcommand{\note}[1]{\par\smallskip{\small\color{gray!60!black}\textit{[#1]}}\par\smallskip}

% A marginal note of hers, distinct from the body of the text.
\newcommand{\marginal}[1]{\par\smallskip\noindent\hspace{1em}%
  {\small\color{gray!60!black}#1}\par\smallskip}

% --- Title ------------------------------------------------------------------
% The title block is French, like the documents themselves.

\AddToHook{shipout/background}{%
  \ifx\grothWatermark\empty\else
    \begin{tikzpicture}[remember picture, overlay]
      \node[rotate=52, scale=3.4, align=center, text=gray!18,
            font=\sffamily\bfseries]
        at (current page.center) {\grothWatermark};
    \end{tikzpicture}%
  \fi}

\AtBeginDocument{%
  \begin{center}
    {\large\bfseries Manuscrits de mathématiciens (Gallica) ---
      \ifx\grothShelfmark\empty\grothFolder\else\grothShelfmark\fi}\\[2pt]
    {\itshape\grothFoldertitle}\\[4pt]
    {\small vues \grothFirst--\grothLast
      \ifx\grothBatch\empty\else\ (lot \grothBatch)\fi}\\[2pt]
    \ifx\grothDating\empty\else
      {\small\color{gray!60!black}Datation du catalogue : \grothDating}\\[2pt]
    \fi
    \ifx\grothWatermark\empty\else
      {\small\bfseries\color{red!45!gray}\grothWatermark}\\[2pt]
    \fi
    {\footnotesize\color{gray!60!black}Transcription établie d'après les images IIIF de Gallica.
     Source gallica.bnf.fr / Bibliothèque nationale de France.\\
     Les lectures incertaines sont soulignées, les illisibles notées [\,\ldots\,],
     les ajouts entre crochets ; \textsf{v.} numérote les vues de Gallica,
     \textsf{f.} la foliotation au crayon de la BnF.}
  \end{center}
  \vspace{1em}}

\endinput


\folder{fr-9115}
\shelfmark{Français 9115}
\batch{36}
\pages{701}{720}
\dating{1801-1900}
\watermark{Lecture automatique\\non vérifiée}
\foldertitle{Papiers de Sophie GERMAIN. — Recueil de dissertations et problèmes mathématiques et physiques.}

\begin{document}

\note{Numérotation des feuillets. En haut à droite, à l'encre, une série d'un nombre par feuillet, sur la face qui le porte : 345 (vue 702, recto blanc), 346 (704), 347 (706), 348 (708), 349 (710), 350 (712), 351 (714, blanc), 352 (716), 353 et 354 (718, deux feuillets sur une même feuille de montage), 355 et 356 (720, de même) ; la vue 700, hors de ce lot, porte 344. Sans lacune ; les revers écrits n'en portent pas. Elle occupe la place de la foliotation de la Bibliothèque mais est tracée à la plume : aucun « f. » n'est donc porté. Les vues 702, 703, 705, 707, 711, 714, 715 et 717 sont des faces blanches, où ne se voit que l'écriture de l'autre face par transparence ; elles ne sont pas transcrites.}

\page{701}

\note{Revers écrit d'un petit feuillet monté bas sur sa feuille ; aucun nombre n'y est écrit. Le recto, vue 700, hors de ce lot, porte en haut à droite, à l'encre, 344, et semble blanc à la vignette : ce qui s'y voit en bas à gauche est l'écriture de cette page lue par transparence. Au-dessus du feuillet paraît le bas d'une autre feuille, blanche, où transparaît une écriture à l'envers. La page commence par une phrase entière et s'arrête au milieu d'une autre, au tiers inférieur du feuillet, le reste blanc.}

\uncertain{Si} par le pôle dela sphère on mène un nombre infini
de méridien et qu'on \struck{\uncertain{fait}} trace en même tems un nombre
infini de parallèles la surface hemisphérique sera
partagée en un nombre infini dequadrilatères rectangulaires
les quadrilatères compris entre deux parallèles \uncertain{seront} tous
plus grands que les quadrilatères qui en nombre égal composent
la zone comprise entre l'un des mêmes parallèles et le suivant
\struck{c'est ce q} car \struck{et} mais comme chacun de ces quadrilatères
est supposé infiniment petit la difference entre ceux qui
composent deux zones contiguës doit être négligeable
ala \ill{}. Si on suppose que sur l'un des côtés méridiens
d'un des parallélogrammes \add{qui s'appuyent sur l'equateur} il existe un foyer \struck{de chaleur}
qui envoye une quantité $f$ de chaleur, que
$f\ f'\ f''\ f'''$\&c soient les \struck{quantités de chaleur qui partent
par les côtés méridiens des differens} temperatures de
differens parallelogrammes qui s'appuient sur l'equateur
$a\ a'\ a''\ a'''$\&c les temperatures des parallelogrammes qui
ont \struck{s'appuient} sur le premier parallèle
\note{« Si » : les premières lettres ressemblent à « ti » ; c'est la forme que prend son \emph{s} initial dans « si on suppose », plus bas sur la même page, d'où la lecture. « seront » est récrit sur un autre mot, et l'ensemble se lit aussi « sauront ». Le mot barré après « qu'on » pourrait être « fait » ou « faire ». Après « négligeable », au début de la ligne suivante, « ala » puis un mot de quatre ou cinq lettres qui se lit « poté » et ne donne pas de sens ; il est laissé illisible. « qui s'appuyent sur l'equateur » est écrit au-dessus de la ligne, d'une plume plus fine, au-dessus de « il existe un foyer ». La page passe de « quadrilatères » à « parallélogrammes » pour les mêmes éléments. La phrase s'arrête sur « sur le premier parallèle », « ont » restant sans complément après la rature de « s'appuient » ; le reste du feuillet est blanc.}

\page{704}

\note{En haut à droite, à l'encre : 346. Petit feuillet monté bas sur sa feuille, écrit sur ses deux tiers inférieurs ; il commence par une phrase entière et finit par un trait, au bas du feuillet. Aucune autre numérotation.}

Voici encore une autre demonstration dela propriété du nombre 2 d'être résidu
ou non résidu des nombres premiers suivant leurs formes par rapport a 8. mais
celle-ci ne peut servir pour distinguer les nombres $8K+5$ des nombres \struck{dans
les quels} $8K+1$ dans lesquels $K$ est divisibles par une puissance paire
de 2.

On peut aumoins demontrer que pour les nombres $4K+3$ 2 est \add{non} résidu lorsque
$K$ est pair et \struck{non} résidu lorsque $K$ est impair en effet $K$ étant $=2K'$
$p=4K+3=8K'+2+1$ \quad $-1\equiv 2(4K'+1)$
$-1$ est non résidu ; $4K'+1$ a un nombre pair de facteurs de laforme
$4h-1$ parconséquent, comme tous les facteurs premiers de $4K'+1$ \struck{son} dela forme
$4h+1$ sont résidus et que ceux dela forme $4h-1$ pris avec le signe $-$ sont
egalement résidus le produit $4K'+1$ est résidu \struck{pa} d'où il résulte 2 non
résidu. Au contraire lorsque $K=2K'+1$ $p=8K'+7=2(4K'+3)+1$, $-1\equiv 2(4K'+3)$
$4K'+3$ a un nombre impair de facteurs dela forme $4h-1$ parconséquent c'est $-(4K'+3)$
au lieu de $4K'+3$ qui est résidu et 2 est alors résidu. car $1\equiv -2(4K'+3)$
\note{Le 3 de « $4K+3$ », de « $4K'+3$ » et de « $2(4K'+3)$ » a la longue queue qui le fait lire 9 ; « $p=4K+3=8K'+2+1$ » et l'argument établissent 3. Dans « $8K'+2+1$ », le 8 est récrit sur un autre chiffre, sans doute un 4. Le « non » ajouté au-dessus de « résidu », à la première ligne de l'alinéa, remplace le « non » biffé de la ligne suivante. Pour la forme des facteurs, la lettre est une petite \emph{h} (« $4h-1$ », « $4h+1$ ») là où les nombres étudiés ont une capitale $K$ ; au début de la ligne « $4h-1$ parconséquent » et dans « $4h+1$ sont résidus » la lettre est récrite sur une autre, peut-être $K'$. Le signe de congruence à trois traits est le sien.}

On peut encore démontrer que 2 est résidu pour les nombres $4n+1$ dans lesquels
$n$ est divisible par une puissance impaire de 2

En effet soit $N$ pairement pair et $2n=N$ \quad $4n+1=2N+1=p$, $-1$ est résidu
(mod.~$p$) aussi bien que tous les facteurs premiers de $N$ et parconséquent leur produit
et $N$ lui même puis qu'il est supposé ne contenir qu'une puissance paire de 2
donc a cause de $-1\equiv 2N$ 2 est aussi résidu. mais on voit que cette methode
n'apprend rien pour le cas $p=2^{2m}n+1$ ---
\note{« tous les facteurs premiers de $N$ » : ses facteurs premiers impairs, puisque $N$ est pairement pair ; la page ne le précise pas. « qu'une » est écrit serré, comme « q'uune ». Dans « $p=2^{2m}n+1$ », « $2m$ » est écrit au-dessus du 2. La page finit sur un trait horizontal.}

\page{706}

\note{En haut à droite, à l'encre : 347. En haut à gauche, un « 1 » suivi d'un trait oblique, de la même encre que le texte : peut-être un numéro de feuillet de sa main, dans la manière de la première rédaction des v.~210–216 (le chiffre suivi d'un trait oblique) ; rien sur la page ne dit à quoi il se rapporte. Le titre est suivi d'un trait horizontal. Le feuillet, monté sur sa feuille, est écrit jusqu'au bord inférieur, et la dernière ligne s'y arrête ; son revers, vue 707, est blanc.}

Observations sur les nombres premiers de la forme $4n+1$
considérés comme résidus d'autres nombres premiers demême forme

Si on a un nombre premier $P=4N+1$ tel que $P$ soit résidu de
$p=4n+1$ $-P$ sera egalement résidu (mod~$p$) ; on aura donc
$X^2\equiv -P$ (mod~$p$) ou $X^2+P=Mp$.

Comme il est demontré d'ailleurs que tout nombre premier de la forme
$4n+1$ est la somme de deux quarrés, on a sans aucune nouvelle\struck{s}
supposition\struck{s}, et en observant que quelque soit $M$ il peut toujours
etre considéré comme la somme de quatre quarrés. et mettant
$A^2+B^2$ au lieu de $P$, $a^2+b^2$ au lieu de $p$ et $p^2+q^2+r^2+s^2$ au lieu
de $M$ ; \quad $X^2+A^2+B^2=(a^2+b^2)(p^2+q^2+r^2+s^2)$
\[=(ap\pm bq)^2+(pb\mp qa)^2+(ar\pm bs)^2+(as\mp br)^2\]
On peut supposer $ap=bq$, $pb+qa=A$ \quad $ar\pm bs=B$ \quad $as\mp br=X$
Des deux premières équations il résulte ; d'abord dela première, a cause de $a$
et $b$ premiers entr'eux, $p=bm$, $q=am$, et dela seconde $A=m(b^2+a^2)$
On auroit encore pû supposer $A=ap+bq$, $pb=qa$ d'où il
résulteroit $p=am'$ \quad $q=bm'$ \quad $A=m'(b^2+a^2)$, c'est adire que
$A$ se réduiroit toujours au produit du nombre premier $p$ par un nombre
quelconque. Comme il est indifferent de prendre $r$ et $s$ au lieu de
$p$ et $q$ puisque les quantités que ces lettres representent ne sont assujetties
a aucune condition particulière, on voit que les suppositions que
l'on vient d'examiner sont les seules par lesquelles on puisse égaler
terme a terme les deux membres de l'équation que l'on considère ;
car, si on fesoit $A=pb\pm qa$ et $B=ap\mp qb$, il en résulteroit
$A^2+B^2=(a^2+b^2)(p^2+q^2)$, c'est adire, que, $A^2+B^2$ ne pourroit
être un nombre premier
\note{La lettre $p$ sert à la fois pour le nombre premier $4n+1$ et pour l'un des quatre carrés de $M$ ; la page ne les distingue pas, et on les laisse tels. Les capitales $P$ et $X$ ont un tracé propre, distinct de $p$ et de $x$. « nouvelles suppositions » : le \emph{s} final des deux mots est barré d'une petite croix, ce qui ramène au singulier (voir hand.md, v.~314). Dans « $A^2+B^2$ ne pourroit », les exposants sont écrits au-dessus des lettres, un peu décalés.}

\marginal{c'est adire par ces manieres que en \uncertain{produisent} aucune}
N.ota Je ne vois pas d'ailleurs que l'on puisse affirmer que
le produit de deux nombres qui ne seroient, l'un et l'autre
décomposables en 4 quarrés que d'une seule façon, ne le soit que de $\underline{2}$
\note{« N.ota » : un N suivi de « ota » en exposant, souligné. La note marginale est écrite en quatre lignes serrées dans la marge de gauche, en face des deux dernières lignes, la dernière sur le bord même du feuillet ; un petit signe en forme de point d'exclamation, entre la note et « décomposables », semble la rattacher à cette ligne. Sa dernière ligne se lit « en produi… aucune », la fin du verbe prise dans une surcharge. La page s'arrête sur le « 2 » souligné, au bord inférieur du feuillet ; le revers est blanc et la suite n'est pas sur la vue suivante écrite (v.~708), qui commence un autre texte.}

\page{708}

\note{En haut à droite, à l'encre : 348. Aucune autre numérotation ; pas de numéro de page de sa main. Le feuillet, monté bas sur sa feuille, est écrit d'une main posée et régulière, presque sans rature ; le texte se poursuit au revers, vue 709. Au-dessus du titre transparaît, à l'envers, l'écriture du revers.}

\emph{Théorème}

$\underline{n}$ étant un nombre impair plus grand que l'unité, il est impossible
de satisfaire en nombres entiers a l'équation $2z^{2n}=y^{2n}+z^{2n}$.
\note{Ainsi sur la page : la dernière lettre de l'équation a la forme du $z$ qui la commence, non celle du $x$ de la ligne « z, y et x premiers entr'eux » ; la même équation, quatre lignes plus bas, s'écrit $2z^{2n}=y^{2n}+x^{2n}$. On garde la lettre écrite. « Théorème » et « Démonstration » sont soulignés, et le $n$ initial l'est aussi.}

\emph{Démonstration}

Nous supposerons comme a l'ordinaire, $z$, $y$ et $x$ premiers entr'eux
et nous observerons que tout diviseur de $z^{2n}$ est nécessairement
la somme de deux quarrés.

Cela posé, soit, s'il se peut, $2z^{2n}=y^{2n}+x^{2n}$ et parconséquent
$z^n=P^2+Q^2$, $x^n=P^2-Q^2-2PQ$, $y^n=P^2-Q^2+2PQ$. En prenant $z=z'z''$,
$x^{n-1}-x^{n-2}y+\&c=H$, $x^{n-1}+x^{n-2}y+\&c=F$ \quad $z''^n=x^{2n-2}-x^{2n-4}y^2+\&c$
et $2z'^{2n}=x^2+y^2$ on aura : $z'^n=p^2+q^2$, $x=p^2-q^2-2pq$, $y=p^2-q^2+2pq$
d'où on tire
\[\left.\begin{array}{r}
P^2+Q^2=z''^n(p^2+q^2)\\
P^2-Q^2=H(p^2-q^2)\\
2PQ=2Fpq
\end{array}\right\}\ \ldots\ (1)\]

Il faut encore faire $z''^n=U^2+V^2$ ; la première des équations (1)
donne alors $P^2+Q^2=(Uq\pm Vp)^2+(Up\mp Vq)^2$, $P=Uq\pm Vp$ et $Q=Up\mp Vq$. En
substituant ces valeurs de $P$ et de $Q$ dans les deux dernières des
équations (1) elles deviennent
\[\left.\begin{array}{r}
\{U(p+q)\pm V(p-q)\}\{U(q-p)\pm V(p+q)\}=H(p^2-q^2)\\
2\{(U^2-V^2)pq\pm(p^2-q^2)UV\}=2Fpq
\end{array}\right\}\ \ldots\ (2)\]

On voit par la première que $VU$ doit être divisible par $p^2-q^2$ ; la
seconde montre que le même produit $VU$ doit aussi être divisible par $pq$.

Soit $z''^{n-1}=V'^2+U'^2$, $z''=v^2+u^2$ et parconséquent
\[z''^n=U^2+V^2=(U'u\pm V'v)^2+(U'v\mp V'u)^2,\ U=U'u\pm V'v,\ V=U'v\mp V'u,\]
\[VU=(U'^2-V'^2)vu\pm(v^2-u^2)V'U'.\]

Les valeurs de $VU$ données par l'un et l'autre signes devront également
satisfaire ; ainsi, leur somme $2(U'^2-V'^2)vu$, et leur différence, $2(v^2-u^2)U'V'$,
seront l'une et l'autre divisibles par $pq(p^2-q^2)$.

En supposant, comme on le doit $V'$ et $U'$ premiers entr'eux, $U'^2-V'^2$
et $V'U'$ seront également premiers entr'eux, parconséquent $V'U'$ ne peut
avoir d'autre diviseur commun avec $2(U'^2-V'^2)vu$ que le produit $\underline{vu}$.

Par la même raison $U'^2-V'^2$ ne peut avoir d'autre diviseur commun
\note{Les capitales $U$, $V$ et les minuscules $u$, $v$ sont nettement distinctes sur la page : $U$ bouclé en tête comme un $\upsilon$, $V$ tracé à gros traits, $u$ et $v$ petits. Les accolades des systèmes (1) et (2) sont les siennes, avec les points de suspension devant le numéro. Dans (2), le signe $=$ de la seconde équation est pris dans un pâté. La page s'arrête au bas du feuillet sur « diviseur commun » ; la phrase continue au revers, vue 709.}

\page{709}

\note{Revers écrit du feuillet 348 (vue 708), sans numéro ; il continue la phrase de la vue précédente (« Par la même raison $U'^2-V'^2$ ne peut avoir d'autre diviseur commun »). Même main régulière, sans rature. La page s'arrête aux trois quarts du feuillet sur l'équation finale, suivie d'un petit trait horizontal : la démonstration est complète.}

avec $2(v^2-u^2)U'V'$ que le nombre $v^2-u^2$ ; il faut donc que $2(v^2-u^2)vu$ soit
divisible par $(p^2-q^2)pq$.

On ne peut satisfaire a cette condition qu'a l'aide de l'une ou l'autre
des suppositions : $2(v^2-u^2)vu>pq(p^2-q^2)$, $2(v^2-u^2)vu=0$.

Examinons d'abord la première.

En mettant pour $p$ et $q$ leurs valeurs, $2(v^2-u^2)uv>pq(p^2-q^2)$ donne
$32(v^2-u^2)vu>2(y^2-x^2)$. A cause de $2(y-x)=8pq$ et $y+x=2(p^2-q^2)$, $2(y^2-x^2)$ est $>(x+y)^2$.
La même supposition donne donc aussi $32(v^2-u^2)vu>(x+y)^2$.

Il résulte de l'équation $2z^{2n}=x^{2n}+y^{2n}$ que $z^{2n}=z'^{2n}(v^2+u^2)^{2n}$ est
plus petit qu'$y^{2n}$. $z'^2(v^2+u^2)^2<y^2$ et a plus forte raison $2z'^2(v^2+u^2)vu<y^2$.
Il faut encore observer que $z'$ étant impair et la somme de deux
quarrés on a $z'>4$ ; parconséquent $32(v^2+u^2)vu<2z'^2(v^2+u^2)vu$ et
$32(v^2+u^2)vu<y^2$. On ne peut donc supposer dans aucun cas
$32(v^2-u^2)vu>(y+x)^2$.
\note{« $2(y-x)=8pq$ et $y+x=2(p^2-q^2)$ » : le mot entre les deux égalités est un signe bref, comme un \emph{d} barré, qu'on rend « et ». Les deux inégalités qui précèdent « On ne peut donc » portent $(v^2+u^2)$, avec le signe $+$, là où la supposition examinée porte $(v^2-u^2)$ ; la page est donnée telle quelle. « qu'$y^{2n}$ » : l'apostrophe et le $y$ sont liés.}

Voyons a présent quels sont ceux dans lesquels on peut
satisfaire a l'équation $2(v^2-u^2)vu=0$.

Si on prend $v^2-u^2=0$, les nombres $x$ et $y$ ne pourront être
premiers entr'eux ; c'est cequ'il est aisé de voir a l'aide des
équations (2) car les quantités $F$ et $H$ ayant alors un
commun diviseur, il faudra que $\underline{x}$ et $\underline{y}$ ayent pareillement
aussi un diviseur commun.

Si on veut que l'un des facteurs $\underline{v}$ ou $\underline{u}$, $\underline{v}$, par exemple,
soit nul, on verra de même que $\underline{x}$ et $\underline{y}$ ne pourront être
premiers entr'eux que dans le seul cas ou l'on auroit
en même tems $u=1$. Soit donc $u=1$, $v=0$, on trouvera
$U=1$, $V=0$, $P=p$, $Q=q$, $x^n=x$, $y^n=y$ et enfin $n=1$. Ce cas
est donc le seul dans lequel il soit possible de
satisfaire a l'équation $2z^{2n}=y^{2n}+x^{2n}$.

\page{710}

\note{En haut à droite, à l'encre : 349. Même main régulière que les vues 708–709, dont ce feuillet est la suite : il revient sur « le théorème précédent » et en étend la démonstration. Aucune autre numérotation. Le feuillet est écrit sur ses trois quarts supérieurs ; le bas est blanc, et son revers, vue 711, l'est aussi.}

La même démonstration s'appliqueroit également a l'impossibilité
de l'équation $2z^{2n}=y^{2(n-2d)}+x^{2(n-2d)}$ ; il faudroit seulement
mettre $n-2d$ au lieu de $\underline{n}$ dans les valeurs de $F$ $H$ et $z''^{2n}$ cequi
ne changeroit rien a la forme des équations (1) et (2).

On seroit mené, comme dans le cas du théorème précédent,
a l'une ou l'autre des suppositions $2(v^2-u^2)vu>pq(p^2-q^2)$
ou $2(v^2-u^2)vu=0$. L'impossibilité dela première résulteroit,
à plus forte raison dela condition $z'^{2n}(v^2+u^2)^{\uncertain{8n}}<y^{2n-4d}$.
La seconde donneroit toujours $u=1$, $v=0$, $V=0$, $U=1$, $P=p$, $Q=q$
et parconséquent $x^{n-2d}=x$, $y^{n-2d}=y$ \quad $n-2d=1$.
\note{L'exposant de $(v^2+u^2)$ est écrit petit et serré : il se lit « 8n », peut-être « 3n » ; le raisonnement parallèle de la vue 709 porte $(v^2+u^2)^{2n}$, et l'on attendrait ici $2n$. On donne ce que montre la page, sous réserve. Dans « $x^{n-2d}$ », le $n$ est repassé d'un trait plus gras.}

En mettant $\underline{n}$ au lieu de $n-2d$ et réciproquement $n+2d$ au lieu
de $\underline{n}$ on pourra donc établir le théorème suivant :

Si $\underline{n}$ est un nombre impair plus grand que l'unité, quelque
soit d'ailleurs le nombre $\underline{d}$, il est impossible de satisfaire
en nombres entiers a l'équation $2z^{2(n+\uncertain{2}d)}=x^{2n}+y^{2n}$
\note{Le chiffre devant $d$, dans l'exposant, est tracé d'un trait lourd, comme récrit sur un autre ; la forme est celle d'un 2, et la phrase précédente (« $n+2d$ au lieu de $n$ ») appuie la lecture. Un petit trait horizontal, sous l'équation, sépare cet alinéa du théorème qui suit.}

\emph{Théorème}

Pour toute valeur de $\underline{n}$ plus grande que l'unité il est impossible de
satisfaire en nombres entiers a l'équation $Z^{2n}=X^{2n}+Y^{2n}$.

\emph{Démonstration}

Si l'équation $Z^{2n}=X^{2n}+Y^{2n}$ étoit possible l'un des nombres $X$ et $Y$
seroit pair on auroit donc en prenant $X=2xx'$ \quad
\[\begin{array}{rl}
x&=\ldots\,(x)\\
&=\ldots\,(x)'\\
&=\ldots\,(x)''
\end{array}\]
\[\begin{array}{ll}
Z^n\pm Y^n=2x'^{2n} & \\
Z^n\mp Y^n=2^{2n-1}x^{2n} & Z\mp Y=2^{2n-1}(x)^{2n}\\
Z^n=x'^{2n}+2^{2n-2}x^{2n} & Z-x'^2=2^{2n-2}(x)'^{2n}\\
\pm Y^n=x'^{2n}-2^{2n-2}x^{2n} & x'^2\mp Y=2^{2n-2}(x)''^{2n}
\end{array}\]
et parconséquent $2^{n-1}(x)^{2n}=2^{2n-2}\{(x)'^{2n}+(x)''^{2n}\}$, \quad $2(x)^{2n}=(x)'^{2n}+(x)''^{2n}$
équation dont l'impossibilité vient d'être demontrée pour toute
valeur de $\underline{n}$ plus grande que l'unité.
\note{Dans le théorème et sa démonstration, $Z$, $X$ et $Y$ sont des capitales droites, distinctes des minuscules $x$, $x'$ ; on garde la distinction. Les trois lignes « $x=\ldots(x)$ » sont écrites à droite de « en prenant $X=2xx'$ », les points de suspension étant les siens ; au-dessus des accents de $(x)'$ et de $(x)''$, une petite tache pâle. Les équations sont disposées sur deux colonnes, la colonne de droite commençant à la deuxième ligne de celle de gauche ; on garde cette disposition. Dans « $Z-x'^2=2^{2n-2}(x)'^{2n}$ », le 2 initial est pris dans un pâté. « $2^{n-1}$ » : l'exposant se lit $n-1$, là où la colonne de droite porte $2^{2n-1}$ ; on le laisse tel. La démonstration s'arrête sur « plus grande que l'unité. », aux trois quarts du feuillet.}

\page{712}

\note{En haut à droite, à l'encre : 350. Aucune autre numérotation. Feuillet monté bas sur sa feuille, écrit jusqu'au bord inférieur, d'une main rapide ; le texte continue au revers, vue 713. Il commence par une équation, sans titre ni phrase d'introduction : le début du calcul n'est pas sur ce feuillet, et rien n'indique qu'il se trouve aux vues précédentes.}

\[\frac{d^2v}{dy^2}=\sin^2\theta\sin^2\varphi.\frac{d^2v}{dr^2}+2\frac{\cos^2\theta\sin.\theta\sin^2\varphi}{r}.\frac{d^2v}{dr.d.\sin\theta}+2\frac{\sin.\varphi\cos.\varphi}{r}.\frac{d^2v}{dr\,d\varphi}\]
\[+\frac{\cos^4\theta\sin^2\varphi}{r^2}.\frac{d^2v}{(d.\sin.\theta)^2}+2\frac{\cos^2\theta\cos.\varphi\sin.\varphi}{r^2\sin.\theta}.\frac{d^2v}{d\varphi\,d.\sin\theta}+\frac{\cos^2\varphi}{r^2\sin^2\theta}.\frac{d^2v}{d\varphi^2}\]
\[+\frac{dv}{dr}\left\{\frac{\cos^2\theta\sin^2\varphi}{r}+\frac{\cos^2\varphi}{r}\right\}\]
\[+\frac{dv}{d.\sin.\theta}\left\{\frac{\cos^2\varphi\cos^2\theta}{r^2\sin.\theta}-\frac{\sin.\theta\cos^2\theta\sin^2\varphi}{r^2}-\frac{2\sin.\theta\cos^2\theta\cos^2\uncertain{\varphi}}{r^2}\right\}\]
\[\uncertain{-}\ \frac{dv}{d\varphi}\left\{\frac{\cos.\varphi\sin.\varphi}{r^2\sin^2\theta}+\frac{\cos^2\theta\sin.\varphi\cos.\varphi}{r^2\sin^2\theta}+\frac{\sin.\varphi\cos.\varphi}{r^2}\right\}\]
\note{Les « $\sin^2\varphi$ » des deux premières lignes ont le mot « sin » récrit à l'encre épaisse sur un autre, peut-être « cos ». À la troisième ligne, le dernier terme touche le bord droit du feuillet : après « cos » on voit l'exposant 2 et le début d'une lettre, qu'on offre comme $\varphi$ ; l'accolade fermante est sur le bord. Le signe qui ouvre la quatrième ligne est un gros pâté allongé, là où l'on attend un signe ; on l'offre comme $-$, sous réserve (voir hand.md, v.~238).}

\[\frac{d^2v}{dx^2}+\frac{d^2v}{dy^2}=\sin^2\theta.\frac{d^2v}{dr^2}+2\frac{\cos^2\theta\sin.\theta}{r}.\frac{d^2v}{dr.d.\sin\theta}+\frac{\cos^4\theta}{r^2}.\frac{d^2v}{(d\sin\theta)^2}\]
\[+\frac{1}{r^2\sin^2\theta}.\frac{d^2v}{d\varphi^2}\]
\[+\frac{dv}{dr}\left\{\frac{\cos^2\theta}{r}+\frac{1}{r}\right\}+\frac{dv}{d.\sin.\theta}\left\{\frac{\cos^2\theta}{r^2\sin.\theta}-\frac{\sin.\theta\cos^2\theta}{r^2}-2\frac{\sin\theta\cos^2\theta}{r^2}\right\}\]

\[\frac{d^2v}{dz^2}=\cos^2\theta.\frac{d^2v}{dr^2}-2\frac{\cos^2\theta\sin.\theta}{r}.\frac{d^2v}{dr.d.\sin\theta}+\frac{\sin^2\theta\cos^2\theta}{r^2}.\frac{d^2u}{(d\sin.\theta)^2}\]
\[\frac{\sin^2\theta}{r}.\frac{dv}{dr}+\frac{dv}{d.\sin\theta}.\left\{\frac{\sin\theta\cos^2\theta}{r^2}+\frac{\sin\theta\cos^2\theta}{r^2}\ \uncertain{-}\ \frac{\sin^3\theta}{r^2}\right\}\]
\note{« $d^2u$ » : ainsi sur la page, là où partout ailleurs la fonction est $v$. La seconde ligne de $\frac{d^2v}{dz^2}$ commence sans signe. Entre ses deux derniers termes, un pâté semblable au précédent tient la place du signe ; on l'offre comme $-$.}

N.ota pour l'intelligence de cette équation il faut observer que la
valeur de la differentielle de $\cos\theta$ prise par rapport a $\sin\theta$
est $\frac{-\sin\theta}{\sqrt{1-\sin^2\theta}}$ et que cette valeur doit être multipliée par
$\frac{d.\sin.\theta}{dz}$ cequi donne $\frac{\sin^2\theta}{r}$. On trouve en effaçant cequi se détruit
\[\frac{d^2v}{dx^2}+\frac{d^2v}{dy^2}+\frac{d^2v}{dz^2}=\frac{d^2v}{dr^2}+\frac{2}{r}.\frac{dv}{dr}+\frac{\cos^2\theta}{r^2}.\frac{d^2v}{(d\sin.\theta)^2}+\left\{\frac{\cos^2\theta}{r^2\sin\theta}-\frac{\sin\theta}{r^2}\right\}\frac{dv}{d.\sin.\theta}\]
\[+\frac{1}{r^2\sin^2\theta}.\frac{d^2v}{d\varphi^2}.\]
\note{« N.ota » : un N suivi de « ota » en exposant, souligné. Dans $\sqrt{1-\sin^2\theta}$, le « sin » est récrit à l'encre épaisse.}

Si on fait $r$ infini l'équation appartient au cylindre on a alors
$\theta=0$ parconséquent $\sin\theta=0$ \quad $\cos\theta=1$ \quad parconséquent elle se réduit à
\[\frac{d^2v}{dx^2}+\frac{d^2v}{dy^2}+\frac{d^2v}{dz^2}=\frac{d^2v}{dr^2}+\frac{1}{r^2\sin^2\theta}.\frac{d^2v}{d\varphi^2}+\frac{d^2v}{d(r\sin\theta)^2}+\frac{1}{r\sin\theta}.\frac{dv}{d(r\sin\theta)}\]
cette équation est en effet celle que M.r Poisson a donné en considérant
le cas général c'est-a-dire celui ou le cylindre est échauffé d'une
manière arbitraire. Si on suppose qu'il n'y ait aucun mouvement
dela chaleur dans le sens de l'axe cequi revient à prendre $\underline{r}$ constant
et parconséquent $\frac{d^2v}{dr^2}=0$ l'équation est celle que j'avois trouvé
j'avois remarqué, dès lors, que dans cette supposition l'équation
du cylindre se réduit à n'être que celle dela plaque circulaire
abstraction faite de la déperdition de chaleur due dans la plaque isolée
au contact de l'air. Si au lieu de prendre $\theta=0$ on prend au contraire
$\theta=90$, \struck{le rayon} l'axe du cylindre deviendra le rayon dela plaque
le second membre de l'équation se réduira a $\frac{d^2v}{dr^2}+\frac{2}{r}\frac{dv}{dr}+\frac{1}{r^2\sin^2\theta}.\frac{d^2v}{d\varphi^2}-\frac{\sin\theta}{r^2}\frac{dv}{d\sin\theta}$
\note{« M.r Poisson » : M suivi d'un petit r surélevé ; le nom est écrit vite mais se lit. « celle que j'avois trouvé », en fin de ligne, puis « j'avois remarqué » au début de la suivante : la phrase se répète ainsi sur la page. « au lieu » est repassé à l'encre épaisse. Le dernier terme, $\frac{\sin\theta}{r^2}\frac{dv}{d\sin\theta}$, est écrit en deux étages au bord inférieur droit du feuillet, le second facteur sous le premier. Le texte continue au revers, vue 713.}

\page{713}

\note{Revers écrit du feuillet 350 (vue 712), sans numéro ; il continue le raisonnement du recto, qui s'arrêtait sur l'équation de la plaque circulaire. Même main rapide. La page s'arrête aux trois quarts du feuillet, sur la troisième des quantités de chaleur ; ce qui paraît au-dessous est l'écriture du recto lue par transparence. La suite n'est pas au feuillet suivant, 351 (vue 714), qui est blanc.}

Pour que cette équation soit identique à celle dela plaque
il faut observer que dans le coefficient $\frac{1}{r^2\sin^2\theta}$ de $\frac{d^2v}{d\varphi^2}$
$\sin\theta$ doit être pris $=1$ puis que $r\sin.\theta$ est alors le
rayon entier et que $r$ multiplie en effet l'expression
sin. mais que dans le terme $\frac{\sin\theta}{r^2}.\frac{dv}{d.\sin\theta}$ $r$ ne multiplie
plus $\sin\theta$ ensorte que cette expression peut être censée
n'avoir pas de valeur fixe quoique l'usage l'attribue
toujours au cercle dont le rayon est l'unité.

Si dans ce terme $\frac{\sin.\theta}{r^2}.\frac{dv}{d.\sin\theta}$ on admet que $\sin\theta$
devienne égal au rayon du cercle dont le rayon
est $r$ on aura $\frac{\sin\theta}{r^2}\frac{dv}{d\sin.\theta}=\frac{1}{r}.\frac{dv}{dr}$.
et parconséquent l'équation devient
\[\frac{d^2v}{dx}+\frac{d^2v}{dy^2}+\frac{d^2v}{dz^2}=\frac{d^2v}{dr^2}+\frac{1}{r}\frac{dv}{dr}+\frac{1}{r^2}.\frac{d^2v}{d\varphi^2}\]
qui est en effet celle dela plaque.
\note{Dans « $r\sin.\theta$ est alors le rayon entier », le $r$ est repassé et pris dans un pâté. Après « l'expression », en fin de ligne, un pâté ; la ligne suivante commence par « sin. », que le sens rattache à « l'expression ». « $\frac{d^2v}{dx}$ » : l'exposant du dénominateur manque sur la page.}

On peut arriver a l'équation du cône en considérant
la quantité de chaleur qui pendant l'instant $\underline{dt}$ entre
par les trois faces qui appartiennent a l'élément de ce solide
l'une de ces quantités sera $Kr^2\sin\theta\cos\theta\,d\theta\,d\varphi.\frac{dv}{dr}$
l'autre seroit $Kr\sin\theta\,d\varphi\,dr.\frac{dv}{r\,d\theta}$ si la direction $\frac{dv}{r\,d\theta}$ étoit
perpendiculaire a la face $r\sin\theta\,d\varphi\,dr$ mais comme elle est
oblique et que l'inclinaison dont il s'agit est mesurée
par l'angle $\theta$ cette quantité doit être multipliée par
$\cos.\theta$ et elle devient $Kr\sin.\theta\cos\theta\,d\varphi\,dr.\frac{dv}{r\sin\theta}$
enfin la troisième seroit $Kr\cos\theta\,d\theta\,dr.\frac{dv}{r\sin.\theta\,d\varphi}$.
\note{Le $K$ est une capitale à boucle, proche d'un \emph{k}. Dans « $\frac{dv}{r\sin\theta}$ », à la fin de l'avant-dernière ligne, le dénominateur ne porte pas de différentielle ; ainsi sur la page. La page s'arrête sur ce terme ; la suite n'est pas dans ce lot.}

\page{716}

\note{En haut à droite, à l'encre : 352. Aucune autre numérotation. En tête, au milieu, souligné, le mot « Cone ». Feuillet monté sur sa feuille, couvert de calculs d'une main rapide jusqu'au bord inférieur ; son revers, vue 717, porte d'autres calculs. Le sujet est celui des vues 712–713 (feuillet 350) : le changement de variables des coordonnées $x$, $y$, $z$ aux coordonnées $r$, $\theta$, $\varphi$ dans $\frac{d^2v}{dx^2}+\frac{d^2v}{dy^2}+\frac{d^2v}{dz^2}$, pour le cône ; ce feuillet en donne le détail, que le feuillet 350 reprenait au net. Rien ne dit lequel fut écrit d'abord.}

Cone

\[\frac{dr}{dx}=\sin\theta\cos.\varphi\qquad \frac{d.\sin\theta}{dx}=\frac{\cos^2\theta\cos\varphi}{r}\qquad \frac{d\varphi}{dx}=-\frac{\sin\varphi}{r\sin\theta}\]
\[\frac{dr}{dy}=\sin.\theta\sin.\varphi\qquad \frac{d.\sin.\theta}{dy}=\frac{\cos^2\theta\sin\varphi}{r}\qquad \frac{d\varphi}{dy}=\frac{\cos\varphi}{r\sin\theta}\]
\[\frac{dr}{dz}=\cos.\theta\qquad \frac{d.\sin.\theta}{dz}=-\frac{\sin\theta\cos\theta}{r}\]
\note{Le signe $=$ de la première ligne est tracé à trois traits ; les autres à deux. Dans « $\frac{\cos^2\theta\cos\varphi}{r}$ », le « cos » de $\varphi$ est récrit à l'encre épaisse.}

\[\frac{dv}{dx}=\frac{dv}{dr}\sin.\theta\cos.\varphi+\frac{dv}{d.\sin\theta}.\frac{\cos^2\theta\cos\varphi}{r}-\frac{dv}{d\varphi}.\frac{\sin\varphi}{r\sin\theta}\]
\[\frac{dv}{dy}=\frac{dv}{dr}\sin.\theta\sin.\varphi+\frac{dv}{d.\sin.\theta}\frac{\cos^2\theta\sin.\varphi}{r}+\frac{dv}{d\varphi}.\frac{\cos.\varphi}{r\sin.\theta}\]
\[\frac{dv}{dz}=\frac{dv}{dr}\cos\theta-\frac{dv}{d.\sin\theta}.\frac{\cos\theta\sin\theta}{r}\]

\[\frac{d^2v}{dx^2}=\sin\theta\cos.\varphi\left\{\frac{d^2v}{dr^2}.\frac{dr}{dx}+.\frac{d^2v}{dr\,d.\sin\theta}.\frac{d\sin\theta}{dx}+\frac{d^2v}{dr\,d\varphi}.\frac{d\varphi}{dx}\right\}\]
\[+\frac{dv}{dr}\left\{\cos\varphi\frac{d.\sin\theta}{dx}-\sin\theta.\sin\varphi\frac{d\varphi}{dx}\right\}\]
\[+\frac{\cos^2\theta\cos.\varphi}{r}\left\{\frac{d^2v}{d\sin\theta.dr}.\frac{dr}{dx}+\frac{d^2v}{d.\sin\theta.d\varphi}.\frac{d\varphi}{dx}+\frac{d^2v}{(d.\sin\theta)^2}.\frac{d.\sin\theta}{dx}\right\}\]
\[-\frac{dv}{d.\sin\theta}\left\{\frac{\sin\varphi\cos^2\theta}{r}\frac{d\varphi}{dx}+\frac{\cos^2\theta\,\struck{\sin\theta}}{r^2}\frac{dr}{dx}\right\}\]
\[-\frac{\sin\varphi}{r\sin\theta}\left\{\frac{d^2v}{d\varphi\,dr}.\frac{dr}{dx}+\frac{d^2v}{d\varphi\,d.\sin\theta}.\frac{d.\sin\theta}{dx}+\frac{d^2v}{d\varphi^2}.\frac{d\varphi}{dx}\right\}\]
\[-\frac{dv}{d\varphi}\left\{\frac{\cos\varphi}{r\sin\theta}\frac{d\varphi}{dx}-\frac{\sin\varphi}{r\sin^2\theta}.\frac{d\sin.\theta}{dx}-\frac{\sin\varphi}{r^2\sin\theta}\frac{dr}{dx}\right\}\]
\[=\frac{d^2v}{dr^2}.\sin^2\theta\cos^2\varphi+\frac{d^2v}{d.\sin\theta\,dr}.2\frac{\sin\theta\cos^2\theta\cos^2\varphi}{r}\ \uncertain{-}\ \frac{d^2v}{dr\,d\varphi}.\frac{2\sin\varphi\cos\varphi}{r}\]
\[+\frac{d^2v}{(d\sin\theta)^2}\frac{\cos^4\theta\cos^2\varphi}{r^2}-\frac{d^2v}{d.\sin\theta\,d\varphi}.\frac{2\cos^2\theta\cos\varphi\sin\varphi}{r^2\sin\theta}+\frac{d^2v}{d\varphi^2}\frac{\sin^2\varphi}{r^2\sin^2\theta}\]
\[+\frac{dv}{dr}\left\{\frac{\cos^2\theta\cos^2\varphi}{r}+\frac{\sin^2\varphi}{r}\right\}\]
\[-\frac{dv}{d.\sin\theta}\left\{\struck{\frac{\cos^2\theta\,\ill{}}{r^2}}\ \frac{\cos^2\theta\sin\theta\cos^2\varphi}{r^2}-\frac{\sin^2\varphi\cos^2\theta}{r^2\sin\theta}\right\}\]
\[+\frac{dv}{d\varphi}\left\{\frac{\cos\varphi\sin\varphi}{r^2\sin^2\theta}+\frac{\cos^2\theta\sin\varphi\cos\varphi}{r^2\sin^2\theta}\ \uncertain{+}\ \frac{\sin\varphi\cos\varphi}{r^2}\right\}\]
\note{Le « sin » barré dans le terme en $\frac{dv}{d.\sin\theta}$, au numérateur de $\frac{\cos^2\theta\sin\theta}{r^2}$, est pris dans une surcharge. Au début de l'accolade de la ligne « $+\frac{dv}{dr}\{\ldots\}-\frac{dv}{d.\sin\theta}\{\ldots\}$ », un premier numérateur est barré et noyé dans un gros pâté ; on y lit « cos$^2\theta$ » et rien d'autre ; au-dessus, « sin$^2\varphi$ » barré. Le numérateur gardé, « cos$^2\theta$ sin $\theta$ cos$^2\varphi$ », est écrit sous la ligne, sur le dénominateur $r^2$. Deux pâtés tiennent la place d'un signe : devant $\frac{d^2v}{dr\,d\varphi}$ et devant le dernier terme de la dernière accolade ; on les offre comme $-$ et $+$ d'après les lignes voisines, sous réserve.}

\[\frac{d^2v}{dy^2}=\frac{d^2v}{dr^2}\sin^2\theta\sin^2\varphi+\frac{d^2v}{d\sin\theta.dr}.\frac{2\sin\theta\cos^2\theta\sin^2\varphi}{r}+\frac{d^2v}{dr\,d\varphi}.\frac{2\sin\varphi\cos\varphi}{r}\]
\[+\frac{d^2v}{(d.\sin\theta)^2}\frac{\cos^4\theta\sin^2\varphi}{r^2}+\frac{d^2v}{d.\sin\theta\,dr}2\cos^2\theta\cos\varphi\sin.\varphi+\frac{d^2v}{d\varphi^2}.\frac{\cos^2\varphi}{r^2\sin^2\theta}\]
\[+\frac{dv}{dr}\left\{\frac{\cos^2\theta\sin^2\varphi}{r}+\frac{\cos^2\varphi}{r}\right\}\]
\[+\frac{dv}{d.\sin\theta}\left\{\frac{\cos^2\theta\cos^2\varphi}{r^2\sin\theta}-\frac{\cos^2\theta\sin\theta\sin^2\varphi}{r^2}\right\}\]
\[\uncertain{-}\ \frac{dv}{d\varphi}\left\{\frac{\cos\varphi\sin.\varphi}{r^2\sin^2\theta}+\frac{\cos^2\theta\sin\varphi\cos\varphi}{r^2\sin^2\theta}\ \uncertain{+}\ \frac{\sin\varphi\cos\varphi}{r^2}\right\}\]
\note{« $\frac{d^2v}{dy}$ » au premier membre, sans l'exposant du dénominateur. À la deuxième ligne, le terme en $2\cos^2\theta\cos\varphi\sin.\varphi$ n'a pas de dénominateur, et sa différentielle est écrite « $d.\sin\theta\,dr$ » là où la ligne correspondante de $\frac{d^2v}{dx^2}$ porte $d.\sin\theta\,d\varphi$ ; on garde la page. Deux pâtés encore tiennent la place de signes, devant $\frac{dv}{d\varphi}$ et devant le dernier terme ; même réserve.}

\[\frac{d^2v}{dx^2}+\frac{d^2v}{dy^2}=\sin^2\theta.\frac{d^2v}{dr^2}+2\frac{\sin\theta\cos^2\theta}{r}.\frac{d^2v}{d.\sin\theta\,dr}+\frac{\cos^4\theta}{r^2}.\frac{d^2v}{(d.\sin\theta)^2}+\frac{d^2v}{d\varphi^2}.\frac{1}{r^2\sin^2\theta}\]
\[+\left(\frac{\cos^2\theta}{r}+1\right)\frac{dv}{dr}+\left\{\frac{\cos^2\theta}{r^2\sin\theta}-\frac{\cos^2\theta\sin\theta}{r^2}\right\}\frac{dv}{d\sin\theta}+\struck{\frac{2\sin\varphi\cos\varphi}{r^2\sin^2\theta}.\frac{dv}{d\varphi}}\]
\note{Le dernier terme est barré de trois traits obliques. « $\left(\frac{\cos^2\theta}{r}+1\right)$ » : le $1$ est écrit à la hauteur de la barre de fraction, hors du numérateur, et sans $r$ ; on le garde tel.}

\[\frac{d^2v}{dz^2}=\cos^2\theta\frac{d^2v}{dr^2}-2\frac{\sin\theta\cos^2\theta}{r}.\frac{d^2v}{d.\sin\theta\,dr}+\frac{\sin^2\theta\cos^2\theta}{r^2}\frac{d^2v}{(d\sin\theta)^2}+(\struck{\ill{}}\,\frac{\sin\theta\cos^2\theta}{r^2}\frac{dv}{d.\sin\theta}\]
\[\frac{d^2v}{dx^2}+\frac{d^2v}{dy^2}+\frac{d^2v}{dz^2}=\frac{d^2v}{dr^2}+\left(\frac{\cos^2\theta+1}{r}\right)\frac{dv}{dr}+\frac{\cos^2\theta}{r^2}.\frac{d^2v}{(d\sin\theta)^2}\]
\[+\frac{1}{r^2\sin^2\theta}.\frac{d^2v}{d\varphi^2}+\left\{\frac{\cos^2\theta}{r^2\sin\theta}+\frac{\sin\theta\cos^2\theta}{r^2}\right\}\frac{dv}{d.\sin\theta}\]
\note{Dans la ligne de $\frac{d^2v}{dz^2}$, la parenthèse ouvrante du dernier terme est suivie d'un signe ou d'une lettre récrits à l'encre épaisse, qui ne se lisent pas. Au dernier terme de la somme, le facteur $\frac{dv}{d.\sin\theta}$ est écrit sous l'accolade, au bord droit.}

en fesant varier par rapport a \struck{$d$} $\sin\theta$. $\cos\theta$ et $\cos^2\theta$
Dans $\frac{dv}{dz}$, $\frac{dv}{dx}$, et $\frac{dv}{dy}$ cette équation devient
\[\frac{d^2v}{dx^2}+\frac{d^2v}{dy^2}+\frac{d^2v}{dz^2}=\frac{d^2v}{dr^2}+\frac{2}{r}\frac{dv}{dr}+\frac{\cos^2\theta}{r^2}.\frac{d^2v}{(d.\sin\theta)^2}+\frac{1}{r^2\sin^2\theta}.\frac{d^2v}{d\varphi^2}+\left\{\frac{\cos^2\theta}{r^2\sin\theta}-\frac{\sin\theta}{r^2}\right\}.\frac{dv}{d.\sin\theta}\]
Si on prend $r$ infini cequi est le cas du cylindre on a en même tems $\sin=0$ \quad $\cos=1$ le
second membre devient $\frac{d^2v}{dr^2}+\frac{1}{r^2\sin^2\theta}.\frac{d^2v}{d\varphi^2}+\frac{d^2v}{d(r\sin\theta)^2}+\frac{1}{r\sin\theta}\frac{dv}{d(r\sin\theta)}$
\note{Les deux lignes de prose sont écrites vite, entre les calculs. La première, « en fesant varier par rapport a $\sin\theta$. $\cos\theta$ et $\cos^2\theta$ », semble expliquer le passage de la somme précédente à la suivante ; le « d » devant $\sin\theta$ est barré. « Dans $\frac{dv}{dz}$ » : le numérateur est récrit sur une autre lettre. La page s'arrête au bord inférieur du feuillet, sur la même conclusion qu'à la vue 712 : l'équation du cylindre.}

\page{718}

\note{Deux pièces sur la même feuille de montage. En haut, un feuillet numéroté en haut à droite, à l'encre, 353 ; en bas à droite, un petit feuillet numéroté 354. Le bord supérieur du feuillet 353 dépasse à la vue 714, où son nombre se lit en partie au-dessus du feuillet 351. Aucune autre numérotation.}

\note{Feuillet 353. De ce côté, deux lignes seulement, à l'endroit, en tête ; tout le reste de ce qui se voit sur la page (un tableau et un raisonnement barrés d'une grande croix) est l'écriture du revers, vue 719, lue par transparence, et n'est pas repris ici. Les deux lignes finissent une phrase dont le début n'est ni sur ce feuillet ni à la vue précédente écrite.}

est non résidu d'un nombre premier de la même
forme $4k+3$
\note{« $4k+3$ » : le 4 a la forme d'un L, comme au petit feuillet ci-dessous (« $4k+3$ », « $4K+3$ ») ; le trait oblique qui touche le 3 est celui de la croix du revers, vu par transparence.}

\note{Petit feuillet 354, écrit sur toute sa hauteur d'une main régulière ; la suite est au revers, vue 719, en bas. Le texte s'ouvre sur « Ici, l'auteur », dans la manière des remarques qu'elle écrit en marge d'une lecture : l'auteur n'est pas nommé.}

Ici, l'auteur ne dit rien des nombres de la forme
$4k+3$.

J'observe que l'on peut démontrer que les
nombres $12n+7$ (lesquels font partie de ceux $4K+3$)
sont non résidus pour quelques nombres premiers
moindres que les nombres dont il s'agit.

En effet si on prend pour $a$ le nombre
pair immédiatement plus grand que $\sqrt{\frac{p}{3}}$,
en fesant $p=12n+7$ \quad $3aa-p$ sera
$<$ que $p$ et de la forme $12K+5$.
or en examinant tous les produits des
nombres non divisibles par 12 on voit
qu'un nombre $12K+5$ a pour facteur
un nombre de la forme $12K+1$ et un
de celle $12K+5$. ou est le produit
de $12K+11$ par $12K+7$ dans le premier
cas soit $12K+5=q$ on aura
$3aa\equiv p$ (mod.~$q$) mais pour les nombres
$12K+5$ $+3$ et $-3$ sont non résidus et
parconséquent $p$ est non résidu (mod $q$)
Dans le second soit $q=12K+7$ on a a cause
de $-3$ résidu, mais $+3$ non résidu (mod.~$q$)
$p$ non résidu même modul.
On peut démontrer la même chose pour
les nombres $12K+5$.
\note{« Ici » : le I est une grande boucle, lu d'après le sens. « ne dit » est écrit « nedit ». « quelques nombres » est écrit « quelquesnombres ». « $<$ que » : le signe est un petit Z couché, comme son signe « plus petit que » ; la phrase le dit en toutes lettres au revers (« sera plus petit que $p$ »). « $3aa-p$ » est ainsi écrit ; le revers, pour $p=12n+5$, écrit $p-3aa$. « ou est le produit » : l'alternative est entre deux décompositions de $12K+5$ ; on garde « ou » tel qu'il est écrit. « modul. » est abrégé ainsi. Le signe de congruence est le sien, à trois traits.}

\page{719}

\note{Revers des deux feuillets de la vue 718, sur la même feuille de montage. Le revers du feuillet 353, en haut, est écrit tête-bêche par rapport à la vue : il se lit en retournant l'image de 180°, et il est barré d'une grande croix à la plume. Le revers du petit feuillet 354, en bas, se lit à l'endroit ; le « 354 » qu'on voit à l'envers en tête est celui du recto, lu par transparence. Aucun nombre n'est écrit de ce côté. On donne d'abord le revers du feuillet 354, qui continue la vue 718, puis celui du feuillet 353.}

\note{Revers du petit feuillet 354. Il écrit pour $p=12n+5$ ce que le recto annonçait (« On peut démontrer la même chose pour les nombres $12K+5$ »). Le bas du feuillet est blanc, à l'écriture du recto près, lue par transparence.}

en effet soit $a$ le nombre pair immédiatement
plus petit que $\sqrt{\frac{p}{3}}$ dans lequel $p=12n+5$
$p-3aa$ sera plus petit que $p$ ce nombre
sera de la forme $12K+5$ et on conclura
comme cidessus qu'il doit être le produit
de $12K+5$ par $12K+1$. ou celui de
$12K+11$ par $12K+7$. d'où il
résulte que dans le premier cas $p$ est
non résidu (mod $12K+5$) et dans
le second non résidu (mod $12K+7$)
\note{« immédiatement » est écrit jusqu'au bord du feuillet, la fin serrée. La phrase finit sans point, et un trait horizontal la sépare du bas du feuillet.}

\note{Revers du feuillet 353, lu la vue retournée. Toute la page est barrée d'une grande croix, deux longs traits qui se coupent vers le milieu ; le texte reste lisible dessous et on le donne en entier, sans le marquer comme biffé.}

A l'égard des nombres $120k+1$ il n'y a qu'a prendre un
nombre $10n+2>\sqrt{\frac{5}{2}.120k+1}$
$2(10n+2)^2-5(120k+1)$ sera un nombre dela forme
$40k+11$ mais pour les nombres $8u+3$ (\uncertain{ou 11}) 2 est non résidu
et pour les nombres $20k+11$ 5 est résidu (\uncertain{ou 19})
donc si $40k+11$ est le produit des deux facteurs
$40k+1$, $40k+11$ et que l'on fasse $4k+11=q$
$2(10n+2)^2\equiv 5(120k+1)$ (mod~$q$) a cause de 2 non résidu
le premier membre est non résidu et parconséquent a cause
de 5 résidu $120k+1$ non résidu (mod.~$q$)

1, 3, 7, 9, 11, 13, 17, 19, 21, 23, 27, 29, 31, 33, 37, 39
\[\begin{array}{lll}
29.39\equiv 11 & q=40k+29 & \text{2 non résidu 5 résidu}\\
9.19\equiv 11 & q=40k+19 & \text{2 non résidu 5 résidu}\\
21.31\equiv 11 & q\equiv 40k+21 & \text{2 non résidu 5 résidu}\\
33.27\equiv 11 & q\equiv 40k+27 & \text{2 non résidu 5 résidu}\\
7.13\equiv 11 & q\equiv 40k+13 & \text{2 non résidu 5 résidu}\\
37.23\equiv 11 & q\equiv 40k+37 & \\
3.17\equiv 11 & &
\end{array}\]
\note{Les deux parenthèses de la troisième et de la quatrième ligne se lisent « (lu 11) » et « (lu 19) » : on les offre comme « ou 11 », « ou 19 », d'après le sens, sous réserve ; ce pourraient être des renvois. « $4k+11=q$ » : ainsi sur la page, sans le 0 que portent les « $40k+11$ » qui précèdent. Dans le tableau, le premier facteur de « $7.13$ » est un 7 pris dans un pâté, qui se lit aussi « 07 » ; $7\cdot 13=91\equiv 11$ (mod 40), comme les autres produits de la colonne, appuie 7. La deuxième colonne passe du signe $=$ au signe $\equiv$ à la troisième ligne ; on garde les signes de la page. La ligne « $37.23\equiv 11$ » s'arrête sur $q\equiv 40k+37$, suivie d'un trait, et « $3.17\equiv 11$ » n'a pas de seconde colonne. Au-dessous, à l'envers de cette lecture, les deux lignes du recto, « est non résidu d'un nombre premier de la même forme », se lisent par transparence.}

\page{720}

\note{Deux feuillets sur la même feuille de montage : en haut, un feuillet numéroté en haut à droite, à l'encre, 355 ; plus bas, un feuillet plus grand numéroté 356. Sur le premier, à droite, le cachet rond de la Bibliothèque impériale. Aucune autre numérotation. Même sujet qu'aux vues 718–719 : les nombres premiers $4k+3$, pris positivement ou négativement, non-résidus d'un nombre premier moindre ; le feuillet 355 nomme Gauss et son article 125.}

\note{Feuillet 355, écrit sur ses deux tiers ; le bas est blanc.}

En examinant la manière dont Gauss démontre (N.o 125) que tout nombre premier
$p=4n+1$ pris negativement est non résidu d'un nombre premier moindre que
$p$ on voit que la même demonstration seroit applicable au cas ou $p$ ne seroit
pas premier desorte que l'on peut affirmer que tout nombre dela forme
$\uncertain{4n+1}$ pris négativement est non résidu de quelque nombre premier moindre
que lui et dela forme $4n+3$. On voit dans la feuille ci-jointe que
tout nombre dela forme $4K+3$ pris négativement est non résidu d'un
nombre \struck{de la} premier de la même forme et moindre que lui il résulte
des ces remarque que : Tout nombre impair $2N+1$ pris avec le signe
$-$ est non résidu d'un nombre premier dela forme $4K+3$
et $<$ que $2N+1$
\note{« Gauss » : le G est une grande boucle ouverte, le nom se lit. « (N.o 125) » : N suivi d'un petit o surélevé, puis 125. Le 4 de « $4n+1$ », « $4K+3$ » a la forme d'un L, comme aux vues 718–719. À la cinquième ligne, le premier chiffre de « $4n+1$ » est récrit et le dernier est pris dans une surcharge ; le sens, qui reprend le cas de la première ligne, appuie la lecture. « $2N+1$ » est écrit au-dessus de « impair », sans signe d'insertion ; on le place là où il est écrit. « des ces remarque » : ainsi sur la page. « $<$ » : son signe « plus petit que », un petit Z couché.}

\note{Feuillet 356, écrit sur toute sa hauteur ; au pied, sous un trait, une remarque d'une plume plus épaisse, et deux notes dans la marge de gauche, appelées par une croix.}

Je me propose de démontrer qu'un nombre quelconque de la forme $4K+3$,
pris, tant avec le signe $+$, qu'avec le signe $-$, est non résidu de quelque
nombre premier moindre que lui.

D'abord soit $p=4K+3$ et soit proposé de prouver que $-p$ est non
résidu d'un nombre $<$ que lui. Prenons $2a$ pour le nombre pair immédiatement
$<\sqrt{p}$ \quad $p-4aa$ sera positif et de la forme $4K+3$ et deplus ce nombre
est évidemment $<$ que $p$. il est clair que $p-4aa$ est premier ou divisible
par un nombre premier de la forme $4K+3$. soit $q$ ce nombre on aura
$p\equiv 4aa$ (mod.~$q$) c'est adire $p$ résidu (mod.~$q$) et parconséquent $-p$ non résidu.

Pour les nombres premiers $4K+3$ positivement pris nous distinguerons deux cas
premièrement si $p$ est dela forme $8n+3$. \struck{soit $a$ un nombre quelconque}
soit $2a$ le nombre pair immédiatement \ill{} que $\sqrt{\frac{p}{2}}$. \quad $p-8aa$ sera dela
forme $8n+3$ et moindres que $p$. ce nombre sera ou premier ou le produit
des facteurs $8K+1$ $8K+3$ \struck{et} $8K+7$ $8K+5$ \struck{et} a cause de 2 non résidu
(mod.~$8K+3$) et (mod.~$8K+5$) nommons $q$ le facteur de l'une de ces
formes qui divise $p-8aa$ on aura $p\equiv 2(2a)^2$ (mod.~$q$) d'où il
résulte $p$ non résidu même modul.
\note{« \struck{soit $a$ un nombre quelconque} » : la fin de la ligne est barrée d'un trait continu ; « quelconque » se devine sous le trait. Le signe entre « immédiatement » et « que $\sqrt{\frac{p}{2}}$ » est noyé dans un pâté : on attendrait $<$ ou $>$, et on ne le lit pas. Les deux « et » donnés comme biffés, entre les paires de facteurs, sont couverts chacun d'une tache d'encre ; peut-être seulement surchargés. « moindres » : ainsi, au pluriel.}

Secondement si $p$ est dela forme $8K+7$. soit $a$ un nombre impair $<$ que $\sqrt{\frac{p}{2}}$
$p-2aa$ sera dela forme $8K+5$ or tout nombre $8K+5$ s'il n'est premier
a pour facteur $8K+1$ et $8K+5$, ou $8K+7$ et $8K+3$ nommons donc
$q$ le facteur de l'une des formes $8K+5$ et $8K+3$ qui divise $p-2aa$
et nous concluerons comme ci-dessus $p\equiv 2aa$ (mod.~$q$) \quad $p$ non résidu (mod.~$q$).

La condition que $2a$ soit immédiatement plus petit de $\sqrt{p}$, dans le premier
cas démontré, et \struck{\ill{}} $\sqrt{\frac{p}{2}}$ dans le second, est évidemment superflue$^{+}$
il suffit de prendre $2a<$ que $\sqrt{p}$ et $2a<\sqrt{\frac{p}{2}}$.
\marginal{$+$ laquelle a l'avantage d'être plus \uncertain{aisée} a \uncertain{remplir}}
\note{La note marginale, en quatre lignes serrées dans la marge de gauche en face du pied du feuillet, est appelée par la croix qui suit « superflue » ; « l'avantage » est écrit « l'avantage » d'un seul trait, « aisée » et « remplir » sont de lecture douteuse, le dernier mot souligné. Le mot biffé devant $\sqrt{\frac{p}{2}}$ est repassé d'un gros trait et ne se lit pas ; peut-être un début de « effectivement ». Le $p$ de « $<$ que $\sqrt{p}$ » est à demi effacé. « $8K+7$ et $8K+3$ » : le premier 8 ressemble à un 6.}

Le but de l'auteur étant de passer au théorème fondamental, il s'est borné a demontrer que les nombres$^{+}$
tant positifs que négatifs sont non résidus. Il faut examiner si on ne pourroit pas faire usage dela même \uncertain{pro.} pour en
\marginal{$4k+3^{+}$}
\note{Ces deux lignes, au pied du feuillet sous un trait horizontal, sont d'une plume plus épaisse et d'une écriture plus serrée. La marque qui suit « les nombres », en fin de ligne, est un signe d'appel, et la marge de gauche porte en face, sous la note précédente, « $4k+3$ » suivi du même signe : on lit donc « les nombres $4k+3$ tant positifs que négatifs ». « l'auteur » est Gauss, que nomme le feuillet 355, et son « théorème fondamental » ; c'est une inférence du transcripteur, la ligne ne le nomme pas. La dernière ligne s'arrête au bord inférieur du feuillet sur « pour en », après une abréviation « pro. » dont la fin se perd : la suite n'est pas sur ce feuillet.}

\end{document}
